\section{Subfield degrees over every odd finite field}
\label{sec:oddpowers}
With the notation of Section~\ref{sec:oddpsl}, one can determine the least degree of a spectrally detecting single subfield for every odd prime power $q\ge5$. Let $h(q)$ be the maximum order of a rich subgroup of $\mathrm{PSL}_2(q)$. Except at four small fields, the answer is
\[
 h(q)=
 \begin{cases}
 q-\sqrt q,&q=s^2,\ s\equiv3\pmod4,\\
 q-q^{3/4},&q\text{ is a fourth power},\\
 (q+1)/2,&\text{otherwise}.
 \end{cases}
 \qquad
 \min_{H\ {\rm rich}} [G:H]=\frac{q(q^2-1)}{2h(q)} .
\]
The four exceptions are $h(19)=12$, $h(41)=24$, $h(101)=h(109)=60$. The index formula always uses the actual value of $h(q)$. In the first two cases an affine subgroup attains equality. For example $h(121)=110$, $h(169)=85$, and $h(625)=500$.

\paragraph{Reduction to an additive hyperplane.}
Write $q=p^f$ and $b=(q+1)/2$. The previous section proves $b\le h(q)<q-1$ for every odd $q\ge5$. Dickson's list, in the form of Theorem~A.1 cited there, excludes all exceptional groups from the order interval $(b,q-1)$ when $q\ge121$, since then $b>60$. For smaller $q$ we test these groups below. A proper subgroup of a torus normalizer has order at most $b$. A subfield $\mathrm{PSL}_2(t)$ with $q=t^j$ has order $(t^3-t)/2$. For $j=2$ its order exceeds $q-1$, and for $j\ge3$ it is less than $b$. A subfield $\mathrm{PGL}_2(t)$ requires even $j$. Its order $t^3-t$ exceeds $q-1$ for $j=2$ and is less than $b$ for $j\ge4$.

For the affine types on Dickson's list, a nontrivial normal elementary abelian $p$-subgroup has a unique common fixed point on the projective line: commuting nonidentity unipotents fix the same point. Its normalizer preserves that point, so the whole group lies in a Borel. The prime-to-$p$ subgroups of the affine type are cyclic of order at most $(q-1)/2$. It therefore remains to consider subgroups in a Borel group $B=U\rtimes T_s$, where $U=(\mathbb F_q,+)$ and $T_s=(\mathbb F_q^\times)^2$ acts by multiplication. For $H\le B$ let $A=H\cap U$ and $C$ be the image in $T_s$. If $A=U$ then $|H|\ge q$, and if $A=0$ then $|H|\le(q-1)/2$. When $C=1$, a proper additive subgroup has order at most $q/p$. Otherwise lift a generator of $C$ to an affine map and translate its unique fixed point to zero. This realizes $H=A\rtimes C$. Put $\mathbb F_r=\mathbb F_p(C)$. The space $A$ is $\mathbb F_r$-linear, so $|A|=q/r^m$ for some $m\ge1$. The condition $|H|>b$ gives $|C|/r^m>1/2$, while $|C|\mid r-1$. It follows that $m=1$ and $C=\mathbb F_r^\times$. Thus $A$ is an $\mathbb F_r$-hyperplane. Since $C$ consists of squares in $\mathbb F_q^\times$, $k=[\mathbb F_q:\mathbb F_r]$ must be even. In particular $q\equiv1\pmod4$.

\paragraph{The exceptional constituent test.}
We prove that $A\rtimes\mathbb F_r^\times$ is rich precisely when
\[
              4\mid\frac{q-1}{r-1}.
\]
Identify the additive characters of $U$ with $\mathbb F_q$ via the absolute trace pairing and a nontrivial additive character of $\mathbb F_p$. The annihilator $A^\perp$ is an $\mathbb F_r$-line, by nondegeneracy of the trace pairing. Its nonzero elements form one free $C$-orbit, contained in one of the two $T_s$-orbits of nontrivial additive characters. Whenever a representation contains a weight in this free orbit, it has an $H$-invariant: take any nonzero vector in that weight space and sum its translates by $C$. The summands have distinct $U$-weights and are all $A$-fixed.

Every ordinary principal, ordinary cuspidal, and Steinberg row contains all nontrivial $U$-weights. Indeed their restrictions to $U$ are, respectively, $\mathrm{reg}_U+1$, $\mathrm{reg}_U-1$ and $\mathrm{reg}_U$, by the unipotent column of the ordinary table. The exceptional pair $E_+,E_-$ consists of the two constituents of $\operatorname{Ind}_B^G\delta$, where $\delta$ is the quadratic character of $T_s$ inflated to $B$. Each has a $\delta$-line on restriction to $B$ by Frobenius reciprocity. Their combined restriction to $U$ is $\mathrm{reg}_U+1$. Hence each has exactly one trivial $U$-line. Each one's other weights fill one of the two $T_s$-orbits, since these have size $(q-1)/2$ and the degree of each exceptional row is $(q+1)/2$. The representation meeting $A^\perp\setminus\{0\}$ has an invariant by the free-orbit argument. In the other, the $A$-invariants are exactly its $\delta$-line. This proves richness if and only if $\delta|_C=1$. The kernel of $\delta$ in $\mathbb F_q^\times$ is the subgroup of fourth powers, so this is exactly the stated divisibility condition.

\paragraph{Maximization.}
For each admissible $r$ the affine subgroup has order $q(1-1/r)$, so choose the largest such $r$, equivalently the smallest admissible even $k$ with $q=r^k$. When $r\equiv3\pmod4$, the sum $1+r+\cdots+r^{k-1}$ is divisible by four for every even $k$. When $r\equiv1\pmod4$, divisibility is equivalent to $4\mid k$. Thus $k=2$ works precisely in the first displayed case. If $4\mid f$ then $k=4$ works, and the square root of $q$ is itself an odd square. If $f$ is odd there is no even $k$. Finally, if $f\equiv2\pmod4$ and $k=2$ fails, then $p\equiv1\pmod4$ and every possible even $k\mid f$ is twice an odd number, so none works. Each of the admissible maximizing affine orders exceeds $b$ whenever the corresponding case occurs ($q\ge9$). This proves the displayed baseline, subject only to checking exceptional subgroups of larger order when $q<121$.

\paragraph{The finite exceptional check.}
Write $h_0$ for the baseline in the cases display. For each odd prime power $5\le q<121$, retain only $K\in\{A_4,S_4,A_5\}$ on Dickson's list with $h_0(q)<|K|<q-1$. For odd $q$, $A_4$ always occurs, $S_4$ requires $q\equiv\pm1\pmod8$ and $A_5$ requires $q\equiv0,\pm1\pmod5$. This yields exactly the following table. Each entry gives $\min_{\chi\in\mathrm{Irr}(G)}\dim V_\chi^K$.
\[
\begin{array}{c|rrrrrrrrrr}
q&17&19&31&41&71&79&81&89&101&109\\\hline
K&A_4&A_4&S_4&S_4&A_5&A_5&A_5&A_5&A_5&A_5\\
\min&0&1&0&1&0&0&0&0&1&1
\end{array}
\]
For reproducibility here is the exact calculation behind the ten entries. The counts of nonidentity elements by order $(2,3,4,5)$ are $(3,8,0,0)$, $(9,8,6,0)$, and $(15,20,0,24)$, respectively. For order $n$ prime to $p$ the torus is split if $n\mid(q-1)/2$, nonsplit otherwise. All cyclic subgroups of that order are conjugate in $G$, and each contains its full set of generators in $K$ whenever it is a subgroup of $K$. For torus parameter exponent $i$, the average of its linear character over these generators is $\mu(d)/\varphi(d)$ for $d=n/\gcd(i,n)$. Thus the average on the class of elements of that order in $K$ is twice this number for a matching principal row, its negative double for a matching cuspidal row, half the corresponding value for a matching exceptional row, and zero for an opposite-series row. The trivial and Steinberg values are given in the previous section. Add the degree to the sum of the order counts times their row averages and divide by $|K|$.

The only defining-characteristic entry is $q=81$, $K=A_5$, $n=3$. All twenty such elements are conjugate in $K$, hence belong to one unipotent class of $G$. The two exceptional values there are $(1\pm9)/2$, with the sign order immaterial. The other row values are $1,-1,0$ for ordinary principal, ordinary cuspidal and Steinberg. These calculations are implemented in \texttt{check\_small\_subgroups.py} with rational arithmetic. They prove exactly the four stated overrides. As a separate check, \texttt{check\_affine.py} constructs finite fields with an explicit primitive element and counts the two unipotent types in nine candidate affine subgroups. For $q=625$ the subgroup using $r=5$ fixes dimensions $1,2$ in the exceptional pair, whereas its larger competitor using $r=25$ fixes dimensions $0,1$.

\paragraph{The remaining odd field.}
For $q=3$ the group is $A_4$. Its three linear characters are inflated from $A_4/V_4$, and its remaining irreducible is the degree-three summand of the permutation character on four letters. Their sum takes values $6,2,0$ at the identity, at involutions and at elements of order three, respectively. It is the unique minimizing character by the same positive-mean argument. A subgroup of order two fixes every linear character and one dimension in the remaining representation, since an involution there has trace $-1$. The degree lower bound $6$ rules out larger rich subgroups, so $h(3)=2$ and the optimal single-subfield degree is $6$.
